\section{引言}

从脑活动中恢复语言是脑机接口研究的重要目标。
侵入式接口已经能够实现大词汇文本解码和接近实时的语音合成
\cite{moses2021neuroprosthesis,willett2023speech,card2024speech,wairagkar2025voice}。
EEG 和 MEG 无需手术植入，但信号较弱。
借助预训练语音和语言模型，非侵入式解码已从连续语音片段检索推进到词级解码
\cite{defossez2023decoding,dascoli2025word}，
其中联合建模相邻词的神经表示可以显著改善预测\cite{dascoli2025word}。

联合建模目标词之后的神经活动，会将解码质量与输入等待耦合在一起。
本文在法语 MEG（Pallier2025）和中文 EEG（ChineseEEG-2）两个自然语音聆听数据集上，
考察词级解码需要多大范围的神经上下文，以及后续输入能否在有限等待下修正已有预测。
目标词左右各最多两个词的局部窗口在宏平均 Top-$10$ 上优于完整上下文模型；
据此，我们让目标词在自身神经窗可用后先给出预测，再随少量后续词窗到达逐步更新，更早的输出依次固定。
贡献如下：
\begin{itemize}
    \item 比较局部词窗与完整上下文在法语 MEG 和中文 EEG 两个自然连续聆听数据条件下的词级检索表现，
          考察较小上下文范围是否足以达到良好的候选排序性能。
    \item 提出共享参数的滚动修正策略，在最多两个后续词窗可用时更新历史词位预测，
          并设置固定 K0 表示的重搜策略，分析表示更新与近期文本搜索的共同作用。
    \item 在词级及序列级评价中考察有限输入等待下的预测质量、回修收益和输出稳定性；
          主要报告法语 MEG 的序列解码比较，并在中文 EEG 上检验相同协议的表现。
\end{itemize}
